%!TEX TS-program = xelatex
\documentclass[11pt, a4paper]{article}

\usepackage[fontset=windows]{ctex}
\usepackage{amsmath, amssymb, amsthm}
\usepackage{graphicx}
\usepackage{hyperref}
\usepackage{bookmark}
\usepackage{enumitem}

\newtheorem{problem}{题}
\renewcommand{\proofname}{\textbf{解答}}
\usepackage[margin=1in]{geometry}

\title{算法设计与分析 作业五}
\author{xxxxxxxxxx xxx}
\date{\today}

\begin{document}

\maketitle

\begin{problem}[7.4]
\end{problem}

\begin{proof}
    设 $N=(V,A,c,s,t)$。因为 $N$ 是简单容量网络，所以任意弧 $a\in A$ 的容量均为
    $c(a)=1$，并且任意两点之间至多有一条弧。

    考察残量网络 $N(f)$ 中由原网络的一条弧 $a=(u,v)$ 产生的残量弧。由于
    $f$ 是 $0$-$1$ 可行流，故 $f(a)\in\{0,1\}$。若 $f(a)=0$，则正向残量容量为
    \[
        c_f(u,v)=c(a)-f(a)=1,
    \]
    反向残量容量为 $0$，因此只可能保留一条容量为 $1$ 的正向弧 $(u,v)$。若
    $f(a)=1$，则正向残量容量为 $0$，而反向残量容量为
    \[
        c_f(v,u)=f(a)=1,
    \]
    因此只可能保留一条容量为 $1$ 的反向弧 $(v,u)$。

    于是，每条原弧在 $N(f)$ 中恰好对应为一条容量为 $1$ 的残量弧，方向或者保持不变，
    或者反向。又因为原网络 $N$ 中任意两点之间至多有一条弧，所以这样的方向替换不会
    产生重边。因此 $N(f)$ 中每条弧的容量仍为 $1$，且任意两点之间至多有一条弧，
    故 $N(f)$ 也是简单容量网络。
\end{proof}

\vspace{1em}

\begin{problem}[7.6]
\end{problem}

\begin{proof}
    设原网络为 $N=(V,E,c,S,T)$，其中 $S$ 为所有发点的集合，$T$ 为所有收点的集合。
    构造一个标准最大流网络 $N'=(V',E',c',s',t')$：
    \[
        V'=V\cup\{s',t'\}.
    \]
    对每个发点 $u\in S$，加入一条弧 $(s',u)$；对每个收点 $v\in T$，加入一条弧
    $(v,t')$；原来的弧全部保留。即
    \[
        E'=E\cup\{(s',u):u\in S\}\cup\{(v,t'):v\in T\}.
    \]
    对原弧 $e\in E$ 令 $c'(e)=c(e)$。新加入的弧取足够大的容量，例如
    \[
        M=\sum_{e\in E} c(e),
    \]
    并令 $c'(s',u)=M,\ c'(v,t')=M$。

    下面说明两问题等价。若 $f$ 是原多发点、多收点网络中的一个可行流，则定义
    $N'$ 中的流 $f'$：原弧上令 $f'(e)=f(e)$；对每个 $u\in S$，令
    \[
        f'(s',u)=\sum_{(u,w)\in E} f(u,w)-\sum_{(w,u)\in E} f(w,u),
    \]
    对每个 $v\in T$，令
    \[
        f'(v,t')=\sum_{(w,v)\in E} f(w,v)-\sum_{(v,w)\in E} f(v,w).
    \]
    由于这些新弧容量为 $M$，不会形成额外限制；而中间点仍满足流量守恒，所以 $f'$
    是 $N'$ 中的可行流，且流值等于原流值。

    反过来，若 $f'$ 是 $N'$ 中的可行流，去掉新源点 $s'$、新汇点 $t'$ 及其关联弧，
    并在原弧上令 $f(e)=f'(e)$，则所有非发点、非收点仍满足流量守恒，容量约束也保持
    不变，因此得到原网络中的一个可行流，且流值相同。于是原问题的最大流值等于构造出
    的标准最大流网络 $N'$ 的最大流值。
\end{proof}

\vspace{1em}

\begin{problem}[7.7]
\end{problem}

\begin{proof}
    设原网络为 $N=(V,E,c,s,t)$。对每个中间点
    $u\in V-\{s,t\}$，将它拆成两个点 $u_{\mathrm{in}}$ 和 $u_{\mathrm{out}}$，
    并加入一条弧
    \[
        (u_{\mathrm{in}},u_{\mathrm{out}})
    \]
    其容量为顶点容量 $c(u)$。源点 $s$ 和汇点 $t$ 不拆分。

    对每条原弧 $e=(u,v)\in E$，在新网络中加入一条对应弧
    \[
        (\alpha(u),\beta(v)),
    \]
    其中
    \[
        \alpha(u)=
        \begin{cases}
            u_{\mathrm{out}}, & u\notin\{s,t\},\\
            u, & u\in\{s,t\},
        \end{cases}
        \qquad
        \beta(v)=
        \begin{cases}
            v_{\mathrm{in}}, & v\notin\{s,t\},\\
            v, & v\in\{s,t\}.
        \end{cases}
    \]
    这条对应弧的容量仍设为原来的边容量 $c(e)$。取新网络的源点仍为 $s$，汇点仍为
    $t$，就得到一个只带边容量的标准最大流网络。

    该构造保持最大流值不变。给定原网络中满足边容量和顶点容量约束的可行流 $f$，
    对每条原弧的对应弧赋同样的流量；对每个中间点 $u$，令
    \[
        f'(u_{\mathrm{in}},u_{\mathrm{out}})
        =\sum_{(v,u)\in E} f(v,u).
    \]
    由于原网络满足顶点容量约束，
    \[
        \sum_{(v,u)\in E} f(v,u)\le c(u),
    \]
    所以拆点弧的容量约束成立；而 $u_{\mathrm{in}}$ 和 $u_{\mathrm{out}}$ 处的流量
    守恒正好对应原点 $u$ 的流量守恒。

    反过来，给定新网络中的任一可行流 $f'$，把每条对应原弧的流量还原为原弧流量。
    对任意中间点 $u$，所有进入 $u$ 的原流量都必须先进入 $u_{\mathrm{in}}$，再经过
    拆点弧 $(u_{\mathrm{in}},u_{\mathrm{out}})$ 才能离开，因此
    \[
        \sum_{(v,u)\in E} f(v,u)
        = f'(u_{\mathrm{in}},u_{\mathrm{out}})
        \le c(u).
    \]
    所以顶点容量约束被满足；边容量和流量守恒也由新网络的可行性直接得到。两边可行流
    一一对应且流值相同，故带顶点容量的最大流问题可以转化为标准最大流问题。
\end{proof}

\vspace{1em}

\begin{problem}[7.17]
\end{problem}

\begin{proof}
    设容量--费用网络为 $N=(V,E,c,a,s,t)$，其中 $c(e)$ 是弧容量，$a(e)$ 是单位流量费用。

    \begin{enumerate}[label=(\arabic*)]
        \item 先暂时忽略费用，在网络中求出最大流值 $F^*$。然后在同一个网络上求一个
        流值恰为 $F^*$ 的最小费用可行流，即求
        \[
            \min \sum_{e\in E} a(e)f(e)
        \]
        满足容量约束、除 $s,t$ 外的流量守恒约束，以及 $|f|=F^*$。这就是最小费用最大流。

        具体可用逐次最短路算法实现：从零流开始，在残量网络中令正向残量弧费用为
        $a(e)$，反向残量弧费用为 $-a(e)$；每次找一条从 $s$ 到 $t$ 的最短费用增广路，
        沿其尽量增广，直到总流量达到 $F^*$。若存在负费用弧，可先用 Bellman--Ford
        求初始势函数，此后用带势的 Dijkstra 求最短路。由于最终流值被固定为最大流值，
        得到的可行流在所有最大流中总费用最小。

        \item 设给定费用上界为 $B$。仍先求最大流值 $F^*$。对任意流值 $x$，记
        \[
            C(x)=\text{流值为 }x\text{ 的可行流的最小费用}.
        \]
        若费用均非负，则 $C(x)$ 随 $x$ 单调不减。因此，问题等价于寻找最大的 $x$，
        使得
        \[
            C(x)\le B.
        \]
        当容量为整数时，可以在区间 $[0,F^*]$ 上二分 $x$：每次用最小费用流算法求
        流值为 $x$ 的最小费用 $C(x)$，若 $C(x)\le B$，说明该流值可达到，继续向右；
        否则向左。最后得到的最大可行 $x$ 就是费用不超过 $B$ 的最大流值，相应的最小
        费用流即为所求流。

        也可以直接使用逐次最短路算法：每次选当前残量网络中从 $s$ 到 $t$ 的最短费用
        增广路 $P$。若其单位增广费用为 $d(P)$，当前剩余预算为 $B'$，则当 $d(P)=0$
        时沿 $P$ 按残量容量尽量增广；当 $d(P)>0$ 时沿 $P$ 增广
        \[
            \Delta=\min\{\text{$P$ 的残量容量},\, B'/d(P)\}
        \]
        的流量。若要求整数流，则在 $d(P)>B'$ 时不能再增加一个单位的流量，于是停止。
        这个过程始终保持在当前流值下费用最小，因此在预算用尽或不存在可增广路时得到的流，
        就是费用不超过给定值的最大流。
    \end{enumerate}
\end{proof}

\vspace{1em}

\begin{problem}[7.21]
\end{problem}

\begin{proof}
    \begin{enumerate}[label=(\arabic*)]
        \item 建立一个分季度的最小费用流网络。设源点为 $s$，汇点为 $t$，四个季度对应
        点为
        \[
            v_1,\ v_2,\ v_3,\ v_4.
        \]
        令各季度生产能力、交货量和单位生产成本分别为
        \[
            (u_1,u_2,u_3,u_4)=(25,35,30,20),
        \]
        \[
            (d_1,d_2,d_3,d_4)=(15,20,25,20),
        \]
        \[
            (p_1,p_2,p_3,p_4)=(12.0,11.0,11.5,12.5).
        \]
        每台设备跨一个季度的库存费用为 $h=0.1$。

        网络中的弧分为三类：
        \[
            (s,v_i),\quad i=1,2,3,4,
        \]
        表示第 $i$ 季度生产。其容量为 $u_i$，单位费用为 $p_i$。
        \[
            (v_i,t),\quad i=1,2,3,4,
        \]
        表示第 $i$ 季度交货。为保证每季度交货量恰好满足合同要求，令该弧的下界和上界
        均为 $d_i$，单位费用为 $0$。
        \[
            (v_i,v_{i+1}),\quad i=1,2,3,
        \]
        表示把第 $i$ 季度生产但未交付的设备库存到下一季度。其容量可取
        $+\infty$，单位费用为 $h=0.1$。

        令源点 $s$ 的供给量为
        \[
            \sum_{i=1}^4 d_i=80,
        \]
        汇点 $t$ 的需求量为 $80$，其余节点满足流量守恒。于是，从 $s$ 到 $v_i$ 的流量
        就是第 $i$ 季度的生产量，从 $v_i$ 到 $v_{i+1}$ 的流量就是第 $i$ 季度末库存量。
        该最小费用流模型的目标函数为
        \[
            \min \sum_{i=1}^4 p_i x_i+\sum_{i=1}^3 h I_i,
        \]
        其中 $x_i=f(s,v_i)$，$I_i=f(v_i,v_{i+1})$。该目标正好等于总生产成本加总库存
        费用，因此求得的最小费用流即给出总成本最低的生产计划。
    \end{enumerate}
\end{proof}

\vspace{1em}

\end{document}
